\section{合卷：合纵、连横与游士——言辞怎样走到山口之前}

函谷关外的道路不宽，能同时开来的车马、粮秣和军队却很多。战国外交的难处，正在于每一个使者都可以把“天下”的利害说得极大，而真正要出兵的国家仍须计算本国的城邑、仓廪、边境和继承。后世常以“合纵”“连横”概括这种竞争：前者大意是东方诸国联合以抗秦，后者多指秦与一国或数国分别结交、使其不入合纵。传世文字中又有“合从”“连衡”等写法，未必是一套当时人人遵守、边界分明的两派学说；它们更适合作为理解联盟聚散的后出名称。

这些名称不能把各国写成两排整齐的棋子。魏要守河洛与大梁，韩的都城和道路紧贴秦东出的门户，赵须顾太行与北方边地，齐能以东方财赋和海滨之利选择介入，楚虽有江汉纵深却难免于北面战线，燕则总要量度辽西、蓟与中原的距离。秦的关中腹地与函谷关使它能够择地防守，又让它特别在意近处的韩、魏城邑。盟约所能做的，是暂时改变这些计算的顺序；它不能取消山口、渡口和粮道。

\subsection{把许多恐惧写成一个盟约，仍须把军队带到同一处}

\begin{event}{公元前318年}{五国合纵攻秦}{可靠纪年}{函谷关及其东侧}\eventanchor{WQ-0318-HEZONG}{战国·五国合纵攻秦}
楚、赵、魏、韩、燕合兵向西，试图以共同出兵阻住秦的东进，终未能越过函谷关取得决定性战果。各国忧惧秦的压力，是结盟的结构性原因；但忧惧并不等于目标相同。对韩、魏而言，眼前是近邻城邑和通道的存亡；对赵、楚、燕而言，远道动员还要与本国边境和内部事务争夺粮饷、兵员与注意力。

函谷关把秦军的防御优势压缩到可以经营的关隘，也把东方联军的困难放大为协同问题：谁先到、谁供粮、谁愿意在别国最危险的方向承担损失，皆不能由盟书自动回答。直接后果是这次大规模合纵未能改变秦据关防守的格局；中期影响则是诸国更常在共同抗秦的名义与各自求缓、求地、求援的实际需要之间往返。它不是“纵约无用”的证据，而是联盟若无持续指挥、补给和共同风险分担，便难以把同一份恐惧变成同一支军队。
\end{event}

这次受挫也不等于秦只须闭关自守。秦若要由关中夺取东方城邑，仍得面对联军再起、齐楚牵制和近邻抵抗。前293年的\eventref{WQ-0293-YIQUE}{战国·伊阙之战}之所以沉重，正因伊、洛之间的峡口把韩魏的防线和秦东出的粮道压在同一处；一场胜利能否延续，取决于能否把守军、道路和新得城邑接进既有网络。合纵与连横并非脱离战争的口才，而是各国围绕这些网络不断改写出兵先后的办法。

\subsection{苏秦、张仪：传记保存的是传统，不是一份会谈录}

苏秦、张仪的名字，几乎同“合纵”“连横”不可分离。《史记·苏秦列传》《张仪列传》把二人的游历、说服与成败组织为彼此映照的人物故事；《战国策》又在不同国策中保存了大量以他们或同类说客为主角的长篇言辞。这些材料十分重要，因为它们让人看见战国人如何把利害、名分、恐惧和荣辱说成可供君主选择的语言；但它们主要是战国以后逐层汇集、整理的说辞与叙事，不能当作每次朝见逐字留下的记录。

所谓苏秦“佩六国相印”、张仪一席之间尽解纵约，正有传记文学所需要的整齐与戏剧性。较可把握的，是诸侯廷中确有跨国游说、馈赠、结交和人材争夺，游士也可能借一国的需求进入另一国的权力核心；难以由此坐实的，是某段辞令的原貌、某次转向只由一人造成，或人物履历中每一处前后相接的年月。把苏、张看作两条强大的政治记忆，比把他们看作能单手转动天下的操盘者，更接近材料的限度。

\begin{event}{公元前316年}{秦取巴、蜀}{可靠纪年}{巴蜀与关中西南通道}\eventanchor{WQ-0316-SHU-BA}{战国·秦灭巴蜀}
秦趁蜀国内乱，出兵兼并蜀、巴。传世叙事将张仪与司马错关于先伐韩、先取蜀的争论写得鲜明；它可以显示秦廷确实要在中原近攻与西南经营之间权衡，不能让我们逐句复原殿上的辩论。较稳固的事件核心是，秦以军事和行政扩张取得成都平原及西南通道，使关中以外的粮食、人口与物资有了新的来源。

这件事也校正了“连横”只靠缔约的想象。秦在外交上同远国暂时修好，是否真能向东持续施压，仍取决于可供运输、征发和驻守的腹地。巴蜀的取得不是后来所有胜利的单一原因，却是使能条件之一：它扩大了秦可调用的资源，也使秦廷在面对东方联盟时多了一层回旋余地。
\end{event}

\begin{event}{公元前312年}{丹阳、蓝田的秦楚战事}{可靠纪年}{丹阳、蓝田及秦楚之间}\eventanchor{WQ-0312-QIN-CHU}{战国·秦楚战争}
秦与楚先后在丹阳、蓝田交战，楚军受挫。后世叙事常把楚怀王是否因张仪许地而绝齐，写成一场由巧言骗成的外交悲剧；《史记·楚世家》《张仪列传》与《战国策》保存的说辞，正使这一情节格外完整。其时间、许地数目和对答的细部在不同文本中并不一致，不能视作当日会谈的速记。

楚的选择也不能只用“受骗”解释。它要在齐秦之间分配关系，又要守住西北方向和广阔国内的调度；秦则把外交施压同边境用兵并用。触发因素可以是齐秦关系的转折，决定战局的却还有秦楚接壤处的道路、军粮与动员。直接后果是楚在西部和汉水方向的压力加重；更长的影响，是楚虽仍为大国，却更难把一次结交或绝交化作稳定的安全。
\end{event}

这类故事里的游士，并非只卖弄辞采。他们携带各国的地理、军势、婚姻、宾客和朝廷人事的消息，在信息不完整时替君主排列风险；也会依附权贵、受制于馈赠，或因旧主失势而改投新国。商鞅由魏入秦、吴起往返魏楚、范雎自魏至秦，都说明人才流动可以带来制度与战场经验，而非只有外交口舌。只是外来者能否被任用，仍由君主、相权、宗族和军功集团的关系决定；游说不是悬在国家之上的自由职业。

\subsection{称号、会盟与离心：同盟为何总在成功之后变难}

\begin{event}{公元前288年}{秦齐互称西帝、东帝，旋即废号}{可靠纪年}{秦、齐及列国外交场域}\eventanchor{WQ-0288-DI}{战国·东西帝号}
秦昭襄王与齐湣王一度互称“西帝”“东帝”，不久又废去帝号。王号已在\eventref{WQ-0334-KINGS}{战国·齐魏称王}中被强国相互承认；更高的称号本可抬升外交位置，却也立刻触动其余诸国对单极秩序的戒心。废号不是礼仪无关紧要，恰说明名分必须有军力、盟友和可接受的利益分配相托，才能长期通行。

《史记·秦本纪》《田敬仲完世家》及《战国策》把这一段与苏秦晚年活动等叙事相连。可确认的是帝号一度出现又撤回；至于谁以哪一句话促成或破坏它，仍应放在后出说辞传统中阅读。人事转折不在于某位辩士神奇地改变了天下，而在于秦、齐的君主和大臣发现，夺取超越诸王的名号会让本可分别谈判的邻国更易结成共同的防备。
\end{event}

前284年燕、赵、秦、韩、魏共同攻齐的\eventref{WQ-0284-QI}{战国·燕破齐}，最能显示联盟的两面。各国并不需要对齐怀有同一份怨恨，仍能在齐孤立时一同进军；可是夺取城邑之后，燕军必须独自面对齐地的守城、征收与地方归附，联盟不能替它长期治理。到\eventref{WQ-0279-TIANDAN}{战国·田单复齐}，占领体系裂开，正说明共同破敌与共同守成是两种不同的政治技术。

秦昭襄王用范雎，后来以“远交近攻”概括对外次序，见\eventref{WQ-0266-FANJU}{秦·范雎入秦与远交近攻}。范雎劝昭王先削韩魏赵、不要耗力远伐齐楚，\eventanchor{WQ-0266-FANJU-KING-COUNSELLING}{战国·范雎劝昭王近攻}，说辞经过传世文本的修辞加工，不能将战略归为一人一时的原话。这句断语不该化作一把万能钥匙：远方国家未必总能安抚，近邻也不会因距离近便无力抵抗。它所触及的却是一个可由地理检验的限制——韩、魏、赵的城邑若被秦夺取，较容易同函谷关内外的道路、县邑与粮运相接；齐、楚即使一时交好，也不能替秦驻守太行或伊洛的缺口。外交的效力，因而总受占领与补给的距离裁定。

\begin{event}{公元前241年}{五国再度合纵攻秦}{可靠纪年}{函谷关与关东诸国}\eventanchor{WQ-0241-HEZONG}{战国·五国再度合纵}
赵、楚、魏、韩、燕再度合兵攻秦，仍未能突破函谷关。秦已置东郡并直接压迫中原，列国共同的危机感比前318年更为迫切；但赵须防北方与秦军，楚有南北广域，魏、韩的城邑首当其冲，燕距主要战场最远，各国无法轻易把本国最后的机动兵力交给统一指挥。

联军退却的直接后果，是各国又回到分别守土、分别求援的处境。中期影响并非秦因一场会战便注定统一，而是秦能继续利用对手无法同时互救的间隙，逐一施压。长平后的力量失衡、秦的郡县化前沿与列国各自的继承和财政难题，才共同构成日后统一战争的条件；把前241年写成合纵“最后一次失败”而略过其后选择，会把历史压得过于平直。
\end{event}

\subsection{使者离开之后：盟约的日常成本}

合纵并不是只在函谷关前才发生。前370年魏公子缓在继承斗争中失败、出奔赵，\eventanchor{WQ-0370-WEI-GONGZHONGHUAN-EXILE}{战国·魏公子缓出奔赵}，便使家内的王位竞争成为赵可使用的外交筹码；流亡路线和支持者不能细考，继承危机会跨过国界却是战国常态。桂陵以后齐赵于平陆等地调整边界关系，\eventanchor{WQ-0355-QI-ZHAO-PINGLU}{战国·齐赵平陆会盟}，会地和条款有异说，却可见胜负之后仍须让军队回营、重新处理相邻城邑。前332年前后魏秦使者围绕河西、上郡和人质往还，\eventanchor{WQ-0332-WEI-QIN-DIPLOMATS}{战国·魏秦就河西与人质交涉}，具体说辞常被后人合并；正因如此，不能把“纵”“横”误读成只在大战前才被选择的一次立场。

张仪、公孙衍的仕宦变化也比传说中的两派对决更曲折。张仪在秦廷活动、倡言连横，\eventanchor{WQ-0322-ZHANGYI-QIN}{战国·张仪在秦廷倡连横}，约前328年居相，\eventanchor{WQ-0328-ZHANGYI-QIN-CHANCELLOR}{战国·张仪居秦相}；相位起讫与游说次序在《史记》《战国策》中并不一致。公孙衍由秦入魏，\eventanchor{WQ-0322-GONGSUNYAN-WEI}{战国·公孙衍入魏谋合纵}，再联络诸国，\eventanchor{WQ-0318-GONGSUNYAN}{战国·公孙衍联络五国}。这些人能把相国、宾客与使者网络接进对外决策，却不能独自替诸国供给军粮；五国在关外面对的粮运、统帅和攻守节奏难题，正是\eventanchor{WQ-0318-HEZONG-LOGISTICS}{战国·五国攻秦的协同困境}所概括的那一层现实。

燕国的内乱同样能改变联盟的可用语言。前317年太子平出奔、联络反对子之的力量，\eventanchor{WQ-0317-YAN-PRINCE-PING-FLIGHT}{战国·燕太子平出奔}，齐后来介入便有了可援引的名义；太子所在和齐的承诺并不清楚，不能把名义当作真实的授权文书。前316年秦廷关于先伐蜀还是东进的司马错、张仪之辩，\eventanchor{WQ-0316-SIMACUO-DEBATE}{战国·司马错张仪伐蜀之辩}，亦是后出人物叙事，不能逐句重演。它仍点出一个限制：联盟与游说须同可供经营的腹地竞争，秦取得巴蜀的选择本身改变了它后来面对关东的资源余地。

\subsection{楚与秦之间：承诺怎样变成关隘}

张仪以商於诱楚绝齐，\eventanchor{WQ-0313-SHANGYU-PLEDGE}{战国·张仪以商於诱楚}；楚怀王断齐，\eventanchor{WQ-0313-CHU-BREAKS-QI}{战国·楚怀王断齐}，却未获所许之地，\eventanchor{WQ-0312-CHU-QI-BREAK}{战国·楚绝齐未获商於}。所谓六百里改为六里的故事，\eventanchor{WQ-0313-ZHANGYI-SIX-LI}{战国·商於六百里与六里之争}，以数字与使节言辞制造了极强的戏剧性；它可能是后人把外交失败压缩成一幕机锋，不能当作逐字会谈记录。无论数字如何，边地承诺未能兑现而楚失去齐的即时支援，足以说明联盟一旦断裂，争执就会落到丹阳、蓝田的道路和军队上。

张仪后来劝魏接受连横，\eventanchor{WQ-0310-ZHANGYI-WEI-PERSUASION}{战国·张仪劝魏连横}，约前310年离秦后在魏不久而卒，\eventanchor{WQ-0310-ZHANGYI-DEATH}{战国·张仪离秦后卒}。个人网络可以结束，秦廷将远交近攻、分别处理近邻的倾向保留下来。前304年秦楚会于黄棘，\eventanchor{WQ-0304-QIN-CHU-HUANGJI}{战国·秦楚黄棘会}，又以边地交换暂缓冲突，\eventanchor{WQ-0304-QIN-CHU-HUANGJI-EXCHANGE}{战国·黄棘边地交换}；人质与使节如何安排，\eventanchor{WQ-0304-QIN-CHU-HUANGJI-HOSTAGES}{战国·黄棘会后的质任往来}，均无完整条约。可见的只是双方以暂时交易处理汉水方向的压力，而非真正消除了防备。

前302年魏襄王入秦朝见、秦归蒲阪，\eventanchor{WQ-0302-WEI-XIANG-QIN-COURT-RITUAL}{战国·魏襄王入秦朝见}，仪式的细节简略，却把军事强制转成可见的从属。前301年垂沙战后齐韩魏未能把临时联军变为常设反楚同盟，\eventanchor{WQ-0301-CHUISHA-AFTERMATH}{战国·垂沙后联军离散}；城邑、俘获与后续进军无法由诸国共同分配，\eventanchor{WQ-0301-QI-HAN-WEI-SPOILS}{战国·垂沙战果未能共管}。盟军胜利仍会迅速散开，正是楚与秦各自仍能寻找回旋余地的原因。

\subsection{函谷关的几年：一次成功并不自动变成同盟}

前299年楚怀王应约入武关，\eventanchor{WQ-0299-CHU-HUAI-WUGUAN}{战国·楚怀王入武关}。入关前后的交涉留下很少，秦把会盟地点放在关隘一侧而取得人身主动权，则比任何密谈都更能说明地理如何参与外交。与此同时，孟尝君失秦相位，\eventanchor{WQ-0297-MENGCHANG-DISMISSAL}{战国·孟尝君失秦相}，逃离细节虽被故事化，齐秦关系由此趋向公开对抗。孟尝君返齐后联络韩魏攻秦，\eventanchor{WQ-0298-QI-HAN-WEI-HANGU}{战国·孟尝君联韩魏攻秦}；联军分段供粮、协同前锋，\eventanchor{WQ-0298-HANGU-COALITION-SUPPLY}{战国·攻函谷关的联军粮运}，从现存材料不能还原为完整补给表，却显示关隘将协同难题放大为日常消耗。

前296年秦割地求和，\eventanchor{WQ-0296-QIN-CEDE-HANGU}{战国·秦为函谷危机割地}，三国的持续攻势停止；久役兵员归国、边城待修，\eventanchor{WQ-0296-TRIPLE-ALLIANCE-DEMOBILIZE}{战国·攻秦联军复员}。复员程序无材料可考，不应凭空补出军令；正因联盟必须把人和粮还给各自国家，秦才能在有限让步之后重新逐近施压。前290年韩魏割地时的朝见、质任等从属安排，\eventanchor{WQ-0290-HAN-WEI-HOSTAGES}{战国·韩魏割地后的质任}，人质名单不存，却再次表明和平并非没有权力差异。

秦齐互称东西帝后迅速撤号，\eventanchor{WQ-0288-EMPEROR-TITLE-WITHDRAWAL}{战国·东西帝号撤回}；苏代劝齐湣王废东帝的说辞，\eventanchor{WQ-0288-SUDAI-ABOLISH-DI}{战国·苏代劝齐废东帝}，不能被当作决定来自一人一言。强国发现其余诸王不愿接受新等级，才是这场名号试验迅速破产的较稳解释。齐灭宋之后，赵、楚、魏协调反齐，\eventanchor{WQ-0285-ZHAO-CHU-WEI-ANTI-QI}{战国·赵楚魏协调反齐}，燕得以把复仇嵌入多国各怀盘算的联盟；苏秦在齐被刺的记忆，\eventanchor{WQ-0284-SUQIN-DEATH}{战国·苏秦在齐遇刺}，身份、年代和遗策都有争议，不能为这个联盟补出一位全知的幕后主使。

前279年渑池会上，秦要赵惠文王鼓瑟、蔺相如以秦王击缶相抗，\eventanchor{WQ-0279-MIANCHI-WINE-CEREMONY}{战国·渑池会礼仪相抗}。鼓瑟击缶属列传戏剧化场景，未必保存实际的会盟程序；礼仪被用来争夺谁能羞辱谁，仍是对强弱关系极敏锐的记忆。前273年华阳援韩时赵魏联军难以统一补给与先后，\eventanchor{WQ-0273-HUAYANG-ALLIED-COMMAND}{战国·华阳赵魏援韩协同失灵}，说明这种对等体面不能替代统一的战场指挥。

\subsection{邯郸、河外与秦亡后的回声}

邯郸被围后，平原君遣使向楚魏求援，\eventanchor{WQ-0259-PINGYUANJUN-ENVOYS}{战国·平原君遣使求援}；赴楚合纵的毛遂故事，\eventanchor{WQ-0257-MAO-SUI}{战国·毛遂自荐赴楚}，言辞虽有浓厚文学性，赵楚为救邯郸确认军事合作，\eventanchor{WQ-0258-ZHAO-CHU-ALLIANCE-RATIFICATION}{战国·赵楚救邯郸合作}，则是这一时期更朴素的事实。救赵联军使秦军外线撤退，不能因此说赵楚从此形成无条件同盟；它只在危机最紧的时刻改变了各国付出兵粮的意愿。

前254年秦取魏吴城、韩王朝秦、魏称听令，\eventanchor{WQ-0254-QIN-WUCHENG-HAN-HOMAGE}{战国·秦取吴城、韩王入朝}；秦以朝觐和威慑维持非吞并从属，\eventanchor{WQ-0254-QIN-HEGEMONIC-COURT}{战国·秦以朝觐维持从属}。秦方的措辞不能被误读为韩魏已亡国，却揭示不战而屈也能成为外交效果。前247年信陵君率五国兵使秦军退河外，\eventanchor{WQ-0247-XINLING-FIVE-STATES}{战国·信陵君率五国攻秦}；联军统一指挥的困难，\eventanchor{WQ-0247-FIVE-STATES-COMMAND}{战国·信陵君联军的统一指挥}，以及秦军撤回关中方向，\eventanchor{WQ-0247-QIN-MENGAO-HEWAI-RETREAT}{战国·秦军自河外撤退}，都说明秦也会受挫，而联盟成员仍迅速回到各自防线。

前241年撤军时各国路线、损失与先后难以详考，\eventanchor{WQ-0241-HEZONG-RETURN-ROUTES}{战国·五国合纵撤军的离心}；不能因为同盟未破关便否认其曾迫使秦重新计算近邻压力。更晚的前207年，巨鹿前诸侯军一度观望、项羽胜后转而听命，\eventanchor{QIN-0207-PRINCES-JULU-OBSERVATION}{秦末·巨鹿后诸侯改从项羽}。这一记忆不属于战国纵横家的时代，却为全章留下一面反光：跨国军队可以在共同危险前集结，仍要等到战场、补给和胜负重排各自的风险，才会暂时接受一个指挥中心。

\begin{sourcenote}[把说辞放回它的文本和地面]
\sourcebook{史记·苏秦列传》《张仪列传》《范雎蔡泽列传》与诸国世家}提供人物、王年和外交转折的主要传世线索；\sourcebook{战国策}保留了最丰富的游说辞令；\sourcebook{资治通鉴}有助于把战争和会盟放回前后年代。《华阳国志》补出巴蜀叙事，《项羽本纪》保留秦末诸侯军的后续回声。它们皆非当时连续的外交档案，尤其《战国策》的国别编排、人物归属与长篇对答，须同《史记》、编年材料和现代史源研究互证。关隘、河流、太行通道、城邑距离和可见的战争后果，能帮助检验一种策略在何处可能有效；它们不能替我们证明某一句话曾在殿上原样说出。
\end{sourcenote}

\begin{turningpoint}[制度得失：能结盟，不等于能共同作战]
\term{所得}：游士和盟约使各国能越过本国朝廷的视野，较快识别强邻、交换信息、争取援军，并在力量不均时制造牵制。秦也可借远交、宾客与人才流动争取时间，把外部压力同内部动员衔接起来。

\term{代价}：说辞会把共同利害说得比实际更整齐，君主和大臣却仍要为本国的粮道、城邑、军功与继承负责。合纵一旦要求某国先出兵、久供粮或放弃近处利益，离心便有了具体的理由；连横一旦让一国暂时获利，也会加深他国的不信。战国外交最深的限制不在辞令是否华美，而在国家能力和地理位置能否承受它所许下的行动。
\end{turningpoint}
